\documentclass[12pt]{article}
\usepackage[utf8]{inputenc}
\usepackage[T1]{fontenc}
\usepackage[margin=1in]{geometry}
\usepackage{setspace}
\usepackage{xcolor}
\usepackage[colorlinks=true, urlcolor=blue!60!black, linkcolor=blue!60!black]{hyperref}
\setlength{\parindent}{2em}
\setlength{\parskip}{1.25em}
\renewcommand{\baselinestretch}{1.0}
\begin{document}
\begin{center}
{\LARGE \textit{Poindexter v. Greenhow}, 114 U.S. 270 (1885)\par}
\vspace{0.8em}
{\large Case Brief --- Four-Page Edition\par}
\end{center}
\vspace{1em}

\section{Memory Jogger}

In the lead \textit{Virginia Coupon Case}, a 5--4 Court held that Virginia's tax-receivable bond coupons were valid contracts, not forbidden bills of credit; that the 1882 law refusing them impaired the obligation of contracts; and that a taxpayer could sue the collector who seized his desk because an officer enforcing a void statute is not the State (\href{https://supreme.justia.com/cases/federal/us/114/270/}{114 U.S. 270}).

\section{Detailed Case Facts}

Facing crushing post-Civil War debt, Virginia's Funding Act of March 30, 1871, let creditors exchange pre-war obligations for new long-term bonds with detached semiannual coupons payable to bearer. To make the bonds marketable, the Act declared that the coupons would ``be receivable at and after maturity for all taxes, debts, dues, and demands due the state.''

When the populist Readjuster Party took control of the General Assembly, it sought to stop coupon redemptions that were draining cash tax revenue. The Act of January 26, 1882, reinforced in 1884, forbade collectors to accept coupons for taxes, required payment in gold, silver, treasury notes, or national bank currency, and abolished common-law remedies such as detinue and trespass against collectors; taxpayers had to pay cash under protest and sue for a refund.

Poindexter, a Richmond taxpayer owing \$12.45, tendered matured 1871 coupons plus currency for the fractional balance to Samuel C. Greenhow, the city treasurer. Greenhow refused the coupons and, under Virginia's 1879 distraint statute, seized Poindexter's office desk.

\section{Procedural History}

On April 26, 1883, Poindexter sued Greenhow personally in detinue to recover the desk. A Richmond police justice dismissed for want of jurisdiction. The Hustings Court of Richmond, the highest state court available given the amount in controversy, found jurisdiction but ruled for Greenhow, holding that the 1882 Act validly barred coupon tenders. Poindexter obtained a writ of error, and the Supreme Court heard the case with seven companion \textit{Virginia Coupon Cases}, deciding it on April 20, 1885.

\section{Judicial Votes}

Reversed and remanded, \textbf{5--4}.
\begin{itemize}
\item \textbf{Majority:} Justice Matthews, joined by Justices Field, Harlan, Woods, and Blatchford.
\item \textbf{Concurrences:} None.
\item \textbf{Dissent:} Justice Bradley, joined by Chief Justice Waite and Justices Miller and Gray, reported after \textit{Marye v. Parsons}.
\end{itemize}

\section{Holding}

(1) Tax-receivable coupons are not ``bills of credit,'' because they become receivable only at maturity and are meant to assure redemption of debt, not to circulate as money. (2) The 1871 Act created a binding contract with every coupon holder, and the 1882 and 1884 Acts impaired it in violation of the Contracts Clause and were void in their entirety, including their repeal of common-law remedies. (3) Because a lawful tender equals payment, Greenhow had no valid authority to distrain, so the detinue action ran against him personally, not against the State, and the Eleventh Amendment did not bar it.

\section{Analysis of Opinions}

\textbf{Majority (Matthews): bills of credit.} Under \textit{Craig v. Missouri} and \textit{Briscoe}, a bill of credit is paper issued on a State's credit and designed to circulate as money. Virginia's coupons were detached interest obligations, receivable only at maturity and expected to be retired, not passed hand to hand.

\textbf{Contracts Clause.} Tax receivability was of the essence of the coupon holder's contract. Because a lawful tender equals payment, Poindexter's tender ended Greenhow's authority to distrain. The pay-under-protest remedy did not preserve the contract but destroyed the right to pay in coupons, and the remedy-stripping provisions could not be severed from a scheme designed to defeat the contract.

\textbf{Eleventh Amendment.} Sovereign immunity cannot be evaded by artful pleading, and suits seeking affirmative performance from the treasury are barred under \textit{Louisiana v. Jumel}. But a tort action against an officer who commits an affirmative wrong against property under an unconstitutional statute is different. Distinguishing the \textit{State}, which can act only within constitutional bounds, from its \textit{government}, which can exceed them, Matthews reasoned that an unconstitutional statute is ``as if it were not and never had been,'' leaving Greenhow an ordinary trespasser.

\textbf{Dissent (Bradley).} Not defending repudiation or denying the contract, Bradley called the majority's distinction an unworkable fiction: sovereign immunity bars every unconsented suit designed to coerce performance of state contracts, whoever is named, and letting taxpayers sue revenue officers who follow state law strips the State of control over its debt and revenue. He stressed that immunity applied against a State's own citizens, foreshadowing \textit{Hans v. Louisiana}.

\textbf{Later developments.} \textit{Houston \& Texas Central Railroad v. Texas} reinforced the narrow bills-of-credit definition. \textit{In re Ayers} barred suits merely to compel state contract performance, and \textit{Ex parte Young} cited \textit{Poindexter} in authorizing prospective injunctions against officials enforcing unconstitutional laws.

\section{Examples: Future Applications}

\begin{enumerate}
\item \textbf{Same-Side (Suit Allowed): Seized Delivery Truck.} A revenue officer who refuses matured tax-receivable coupons and seizes a company's truck under a statute repudiating them is answerable personally, as the void statute gives him no authority.
\item \textbf{Same-Side (Not a Bill of Credit): Tax-Anticipation Warrants.} Warrants issued in large denominations, payable in a year and receivable only at maturity are debt, not currency.
\item \textbf{Opposite-Side (Likely Bills of Credit): Small Tax Credit Notes.} Bearer notes in \$5 to \$20 denominations, receivable immediately, resemble \textit{Craig}'s circulating certificates; if void, there is no contract to protect.
\item \textbf{Opposite-Side (Barred): Order to Pay Overdue Coupons.} A suit ordering the treasurer to pay coupons from the treasury seeks the State's own performance, barred under \textit{Jumel} and \textit{In re Ayers}.
\item \textbf{Fence-Sitter: Injunction to Accept Coupons.} A prospective injunction requiring collectors to accept coupons draws on \textit{Ex parte Young} but resembles compelled specific performance, so its characterization likely decides the case.
\end{enumerate}

\section{Critique}

\textbf{The central paradox.} The majority insists that Greenhow is \textit{not} acting for Virginia under the Eleventh Amendment, yet his conduct violates the Contracts Clause, which restricts only ``State'' action, precisely because he enforces Virginia's 1882 Act.

\textbf{An enduring fiction.} The officer is thus a state actor for the merits but a private wrongdoer for immunity, and that paradox passed directly into \textit{Ex parte Young} and remains the defining analytical fiction of federal courts doctrine.

\section{Key Quotations}

\begin{quote}
``The defendant in error is not her officer, her agent, or her representative, in the matter complained of; for he has acted not only without her authority, but contrary to her express commands.''
\end{quote}
\noindent --- Justice Matthews, Opinion of the Court, 114 U.S. at 292--93.

\begin{quote}
``It is well settled by many decisions of this Court that, for the purpose of affecting proceedings to enforce the payment of taxes, a lawful tender of payment is equivalent to actual payment, either being sufficient to deprive the collecting officer of all authority for further action, and making every subsequent step illegal and void.''
\end{quote}
\noindent --- Justice Matthews, Opinion of the Court, 114 U.S. at 282.

\begin{quote}
``They do not become receivable in payment of taxes till they are due, and the design, we are bound to presume, was that they would be paid at maturity. This necessarily excludes the idea that they were intended for circulation at all.''
\end{quote}
\noindent --- Justice Matthews, Opinion of the Court, 114 U.S. at 285.

\end{document}
